% Ad Atticum 11.6 --- headnote
% ad-atticum-11-06, 28 November 48 BC, Brundisium

\textit{Cicero to Atticus, written from Brundisium on the
fourth day before the Kalends of December 48 BC --- 28
November (the manuscript dateline: \textit{Scr.\ Brundisi iv
K.\ Dec.\ a.\ 706 (48)}, and the date is repeated at the
foot of the letter itself). Three weeks have passed since
11.5; Cicero is still pinned at Brundisium and still
defending, to Atticus and to himself, the decision to leave
Pompey's camp. The letter is more substantial than the
previous one, more articulate, and considerably darker. It
opens with the grief that doubles when shared (\S1): Atticus
approves the decision, says the men who matter approve it
too --- but if Cicero really believed that, ``I should feel
the pain less.''}

\textit{The centre of the letter (\S2) is the harshest
formal apologia for Pharsalus-side defection that Cicero
ever puts on paper. The Pompeian camp, he insists, had become
something foreign to the republic: cruelty so unbounded, and
so close a coupling with barbarian peoples, that a proscription
had been ``sketched out not by name but by class,'' with the
property of every senator who had stayed in Italy already
counted as booty of the coming victory. He does not regret
the will, only the plan: better to have sat down in some
town and waited to be summoned than to lie at Brundisium ---
exposed on every side --- and unable now to move closer
without lictors he cannot lawfully shed. The section ends in
two short, manuscript-corrupt clauses preserved here as
\dag\ markers, the second of which has resisted convincing
emendation; the sense given is the most natural reading of
the corrupted text. \S3 turns to operational matters: he has
written to Oppius and Balbus to ask in what manner he ought
to come nearer to Rome, and floats the idea that Trebonius,
Pansa, and others write to Caesar endorsing his action as
taken on the recommendation of Caesar's own circle. \S4 is
the single private cry of the letter --- Tullia is ill, and
the illness is undoing him. \S5 is the dignified, unsurprised
notice of Pompey's death (``About Pompey's end I never had a
doubt''), with one of the gentlest valedictions Cicero ever
gives a dead political opponent. \S6 collects what news there
is from elsewhere: Fannius has been talking dangerously about
Atticus's having stayed in Italy; Lentulus had already
parcelled out the spoils in his head (Hortensius's house,
Caesar's gardens, Baiae); the two sides, he notes drily,
are not very different except that the other was without
limit. He hears Quintus has gone to Asia to plead his case
and has no news of his nephew. The letter closes with a
request for letters back as soon as possible, and the
Brundisium dateline.}
